\documentclass[aspectratio=169]{beamer}
\usepackage{appendixnumberbeamer}
\usepackage{mathtools}
\usetheme{metropolis}           % Use metropolis theme
\definecolor{fhggreen}{RGB}{23,156,125}
\let\oldemph\textbf
\renewcommand{\textbf}[1]{{\color{mLightBrown}\oldemph{#1}}}

\usepackage{blkarray}
\usepackage{amsmath}
\usepackage{multirow}
\title{
\includegraphics[width=0.2\textwidth]{gns-logo.png}\\\small{The GNU Name System}}
\date{2023-09-27}
\author{Bernd Fix, Christian Grothoff, \textbf{Martin Schanzenbach}}
%\institute{\includegraphics[width=.25\textwidth]{aisec_logo.pdf}}

\graphicspath{{figures/}}

\begin{document}
\metroset{block=fill,sectionpage=progressbar,numbering=counter}
\maketitle

\begin{frame}{Directories / DNS}
  $$
    \mathrlap{\overbrace{\phantom{\text{www}}}^{\text{Label}}}
    \text{www}
    \mathrlap{\underbrace{{\color{fhggreen}\phantom{\text{.example.com}}}}_{\text{Namespace}}}
    {\color{fhggreen}\text{.example.com}}
  $$
\end{frame}

\begin{frame}{The .alt TLD}
  Whats wrong with DNS? See RFC 8324\footnote{DNS Privacy, Authorization, Special Uses, Encoding, Characters,
          Matching, and Root Structure: Time for Another Look?}:
  \begin{itemize}
    \item No \textbf{query privacy}.
    \item A \textbf{single hierarchy with a centrally controlled root}.
    \item Requires management/maintenance of \textbf{root servers}.
    \item etc\ldots
  \end{itemize}
  DNSSEC and other ``patches'' do not or in adequately address the issues: ``[the existing solutions for DNS are] security patches rather than designed-in
   security or privacy mechanisms''.
\end{frame}

\begin{frame}{Directories / GNS}
  $$
    \mathrlap{\overbrace{\phantom{\text{www}}}^{\text{Label}}}
    \text{www}
    \mathrlap{\underbrace{{\color{fhggreen}\phantom{\text{.myzone.gns.alt}}}}_{\text{Namespace}}}
    {\color{fhggreen}\text{.myzone.gns.alt}}
  $$
\end{frame}

\begin{frame}{The .alt TLD}
  Why ``.gns.alt''?
  \begin{itemize}
    \item RFC9476: ``The .alt Special-Use Top-Level Domain'' defines the TLD to be used for alternative (from the point of view of DNS) name systems.
    \item RFC9476 does \textbf{not} define a registry for ``.alt''-subdomains.
    \item We manage a ``.alt'' registry at \url{https://gana.gnunet.org}~\footnote{If you ever need a registry for your protocol feel free to approach us!} which already includes a code point for ``.gns.alt''.
    \item To prevent shadowing of DNS names, it is recommended to use the ``.gns.alt'' suffix.
    \item Sometimes (e.g. censorship-overrides) you may not want to do that.
  \end{itemize}
\end{frame}

\begin{frame}{The GNU Name System}
  \begin{itemize}
    \item Namespaces are created and uniquely identified using \textbf{public zone keys}.
    \item \textbf{Records} are grouped by \textbf{label}, encrypted, signed, and published in a key-value store (usually, a DHT\footnote{\url{https://datatracker.ietf.org/doc/draft-schanzen-r5n/}}).
    \item Supported zone types and crypto (for now):
      \begin{itemize}
        \item PKEY: ECDSA+CTR-AES-256
        \item EDKEY: EdDSA+XSalsa20-Poly1305
      \end{itemize}
  \end{itemize}
  \centering
  \includegraphics[width=1\textwidth]{GNS-BlockCreation-0}
\end{frame}

\begin{frame}{Zone management}
  \includegraphics[height=0.9\textheight]{deleg0.pdf}
\end{frame}


\begin{frame}{Name resolution}
  \includegraphics[height=0.9\textheight]{deleg1.pdf}
\end{frame}


\begin{frame}{Name resolution}
  \includegraphics[height=0.9\textheight]{deleg2.pdf}
\end{frame}


\begin{frame}{How do we bootstrap the top-level zones?}
  \includegraphics[height=0.9\textheight]{deleg3.pdf}
\end{frame}


\begin{frame}[fragile]{The Start Zone}
  ``Hyper-hyper local root'' concept we call the \textbf{Start Zone}:
  \begin{itemize}
    \item Start Zone contains so-called \textbf{suffix-to-zone}-mappings.
    \item Implementation ships with an \emph{initial} Start Zone configuration.
    \item Start Zone is configurable \emph{locally} at \emph{each} endpoint.
    \item User override/extension of mappings at top-level or subdomain-level for:
      \begin{itemize}
        \item Circumvent censorship if necessary.
        \item Private networks.
      \end{itemize}
  \end{itemize}
\end{frame}

\begin{frame}[fragile]{The Start Zone}
  Example suffix-to-zone mappings:
  \begin{small}
  \begin{verbatim}
# Some TLDs
.com = 000G001MF6DVMZZ4Y8XRZQDXM1PB3D3VGEK29ZHXBA57EPSNW1QBPKT8J0
.myzone.gns.alt = 000G007FKSA876G6SNDF8VA7YK1DJE96RPPBHRT2X55Q13M2T4YKNYT3DG
# Some subdomain overrides
.gnu.org = 000G001223Q8ZJZBSK6XT2DWV6PE5B1W436D2NB7ZBR9XSXT7TFJHCDB24
.gnunet.gns.alt = 000G0047M3HN599H57MPXZK4VB59SWK4M9NRD68E1JQFY3RWAHDMKAPN30
  \end{verbatim}
  \end{small}
\end{frame}


\begin{frame}{Possible Governance Models}
  \begin{itemize}
    \item Non-profit organization.
    \item Multi-stakeholder model: Board, supporting organizations, \ldots
    \item Examples for possible stakeholders:
      \begin{itemize}
        \item Software and OS Distributors
        \item Browser vendors
        \item Governments
      \end{itemize}
    \item Funding options:
      \begin{itemize}
        \item Applications for new top-level domains.
        \item Registrations of new top-level domains.
        \item \ldots
      \end{itemize}
  \end{itemize}
\end{frame}

\begin{frame}{Hiding information inside GNS}
  \begin{itemize}
    \item GNS's crypto allows you to hide resource records.
    \item It requires either
      \begin{itemize}
        \item the use of a label with sufficient entropy (a shared secret) or
        \item the use of a secret zone.
      \end{itemize}
  \end{itemize}
  $$
    \mathrlap{\overbrace{\phantom{\text{ohcoxaiShaingahd}}}^{\text{Secret label}}}
    \text{ohcoxaiehaingahd}
    \mathrlap{\underbrace{{\color{fhggreen}\phantom{\text{.my.zone}}}}_{\text{Namespace}}}
    {\color{fhggreen}\text{.my.zone}}
  $$
\end{frame}


\begin{frame}{Encrypt}
  \centering
  \includegraphics[width=1\textwidth]{GNS-BlockCreation-1}
\end{frame}

\begin{frame}{Sign}
  \centering
  \includegraphics[height=0.9\textheight]{GNS-BlockCreation-2}
\end{frame}

\begin{frame}{Derive}
  \centering
  \includegraphics[height=0.9\textheight]{GNS-BlockCreation-3}
\end{frame}

\begin{frame}{Combine and publish}
  \centering
  \includegraphics[width=1\textwidth]{GNS-BlockCreation-4}
\end{frame}

\begin{frame}{Query}
  \centering
  \includegraphics[height=0.9\textheight]{GNS-BlockValidation-0}
\end{frame}

\begin{frame}{Retrieve}
  \centering
  \includegraphics[height=0.9\textheight]{GNS-BlockValidation-1}
\end{frame}

\begin{frame}{Verify}
  \centering
  \includegraphics[height=0.9\textheight]{GNS-BlockValidation-2}
\end{frame}

\begin{frame}{Decrypt}
  \centering
  \includegraphics[height=0.9\textheight]{GNS-BlockValidation-3}
\end{frame}

\begin{frame}{Decrypt}
  \centering
  \includegraphics[height=0.9\textheight]{GNS-BlockValidation-4}
\end{frame}

\begin{frame}{Status}
  \begin{itemize}
    \item Specification efforts:
      \begin{itemize}
        \item \url{https://datatracker.ietf.org/doc/draft-schanzen-gns/} -- Will become RFC soon (TM).
        \item \url{https://datatracker.ietf.org/doc/draft-schanzen-r5n/} -- Is being worked on.
      \end{itemize}
    \item Reference implementation in C (part of GNUnet), alternative implementation in Go.
    \item Currently funded project to develop and host a GNS zone registrar service and to mirror some (large) DNS zones funded through NLnet / NGI Zero Entrust.
    \item Current and future research:
      \begin{itemize}
        \item PQ-secure key blinding.
        \item Sharing identity information via GNS (re:claimID).
      \end{itemize}
  \end{itemize}
\end{frame}


\begin{frame}
  \begin{center}
    Questions?\\
    \vspace{2cm}
    \url{https://gnunet.org}\\
    \vspace{1em}
    {\tiny
    \texttt{schanzen@gnu.org}\\
    \texttt{3D11~063C~10F9~8D14~BD24~D147~0B09~98EF~86F5~9B6A}\\
    }
  \end{center}
\end{frame}

\end{document}
